\documentclass[10pt,a4paper,twocolumn]{article}

% ---------------------------------------------------------------------------
% Packages
% ---------------------------------------------------------------------------
\usepackage[a4paper,margin=1.5cm,columnsep=0.6cm]{geometry}
\usepackage[T1]{fontenc}
\usepackage{newtxtext,newtxmath}
\usepackage{microtype}
\usepackage{graphicx}
\usepackage{booktabs}
\usepackage{siunitx}
\usepackage{xurl}  % break long URLs in the references
\usepackage{xcolor}
\usepackage[font=small,labelfont=bf,skip=4pt]{caption}
\usepackage[compact]{titlesec}
\usepackage[colorlinks=true,linkcolor=blue!50!black,citecolor=blue!50!black,urlcolor=blue!50!black]{hyperref}

\graphicspath{{figures/}}

% Tight journal-style headings
\titleformat{\section}{\normalfont\large\bfseries}{\thesection}{0.5em}{}
\titleformat{\subsection}{\normalfont\normalsize\bfseries}{\thesubsection}{0.5em}{}
\titlespacing*{\section}{0pt}{1.0ex plus 0.5ex minus 0.2ex}{0.5ex}
\titlespacing*{\subsection}{0pt}{0.8ex plus 0.3ex minus 0.2ex}{0.3ex}

\setlength{\parskip}{0pt}
\setlength{\textfloatsep}{8pt plus 2pt minus 2pt}
\setlength{\floatsep}{6pt plus 2pt minus 2pt}

% Units
\DeclareSIUnit\year{yr}
% "+0.62 ± 0.33 dB/yr", matching the figures (siunitx 2 defaults: "0.62(33) dB yr^-1")
\sisetup{separate-uncertainty=true, retain-explicit-plus=true, per-mode=symbol}

% Draft helpers: \todo{...} shows red text; remove before submission
\newcommand{\todo}[1]{\textcolor{red}{[#1]}}

% Include a figure if it exists, otherwise draw a labelled placeholder box
\newcommand{\figorbox}[2]{%
  \IfFileExists{figures/#1}%
    {\includegraphics[width=\columnwidth]{#1}}%
    {\fbox{\parbox[c][#2][c]{0.95\columnwidth}{\centering\small\ttfamily figures/\detokenize{#1}}}}%
}

% ---------------------------------------------------------------------------
% Title block
% ---------------------------------------------------------------------------
\title{\vspace{-1.2cm}\Large\bfseries Sentinel-1 C-band backscatter trends on Bodmin Moor peatland, 2023--2024}
\author{Joshua Prettyman}
\date{}

\begin{document}

\twocolumn[
  \begin{@twocolumnfalse}
    \maketitle
    \vspace{-0.8cm}
    \begin{abstract}
      \noindent
      We investigate how Sentinel-1 VV backscatter over peatland on Bodmin
      Moor, Cornwall, varied between January 2023 and December 2024. Over the
      moor as a whole the backscatter increased by
      \SI{+0.62 \pm 0.33}{\deci\bel\per\year}, but this appears to be mainly
      the recovery from a drought in June 2023, and local trends relative to
      the moor are weak. At a site where restoration began in March 2023, the
      backscatter had a smaller annual cycle and a smaller response to the
      drought than a nearby control area, which may indicate a more stable
      water table, although with two years of data we cannot separate this
      from an existing difference between the sites.
    \end{abstract}
    \vspace{0.4cm}
  \end{@twocolumnfalse}
]

% ---------------------------------------------------------------------------
\section{Introduction}
% ---------------------------------------------------------------------------
Satellite radar has been used to monitor the condition of peatlands, both through the intensity of the backscattered signal, which responds to surface wetness and vegetation \cite{toca2023potential,asmuss2019potential,hrysiewicz2023sar}, and through InSAR measurements of surface motion, which responds to the water and gas content of the peat \cite{marshall2021naturescot,bradley2022identification}.
Bodmin Moor, Cornwall, is one of the areas covered by the South West Peatland Partnership \cite{swpp}, and several peatland restoration projects on the moor began in 2022 and 2023 (Section~\ref{sec:study}).

In this report we investigate how Sentinel-1 VV backscatter over peatland on Bodmin Moor varies between January 2023 and December 2024, and where on the moor it shows the clearest increase or decrease.
We then compare a site where restoration began in March 2023 with a nearby control area.

% ---------------------------------------------------------------------------
\section{Data and methods}
% ---------------------------------------------------------------------------

\subsection{Data and acquisition filters}

We use Sentinel-1 GRD data from Earth Engine \cite{gorelick2017}, which are pre-processed with the \emph{Sentinel-1 Toolbox} (thermal noise removal, radiometric calibration and terrain correction) and converted to decibels \cite{eeS1grd}.
Sentinel-1 records both VV and VH polarisations, and we use VV, which in drained grassland peat correlated more strongly with water-table depth than VH \cite{asmuss2019potential}.
Three relative orbits of Sentinel-1A pass over our study area, a bounding box of \ang{4.66}--\ang{4.55}\,W, \ang{50.55}--\ang{50.63}\,N, each giving one acquisition every 12 days (Table~\ref{tab:orbits}).

\begin{table}[h!]
  \centering
  \small
  \caption{Sentinel-1A relative orbits covering the study area, January 2023 to
  December 2024 (IW mode, GRD, VV, \SI{10}{\metre} pixels). Incidence angle:
  mean (range) over the study area. Amplitude: of the annual cycle fitted
  (Eq.~\ref{eq:model}) to the median VV over the masked peat.}
  \label{tab:orbits}
  \setlength{\tabcolsep}{4pt}
  \begin{tabular}{@{}rlcccc@{}}
    \toprule
    Orbit & Pass  & Time  & Dates & Incidence        & Amplitude \\
          &       & (UTC) &       & angle (\si{\degree}) & (dB)      \\
    \midrule
    1     & Asc.  & 18:14 & 58    & 44.8 (44.5--45.0) & 0.92 \\
    52    & Desc. & 06:32 & 61    & 37.7 (37.4--38.0) & 1.16 \\
    103   & Asc.  & 18:06 & 60    & 36.5 (36.2--36.8) & 1.45 \\
    \bottomrule
  \end{tabular}
\end{table}

We use orbit 103 for all figures since it has the smallest incidence angle, so the signal passes through less vegetation, and the largest annual cycle over the moor, and it passes in the evening, avoiding the dew and frost which may affect orbit 52 before sunrise.
We also use orbits 1 and 52 for comparison. 

\subsection{Study areas}
\label{sec:study}

We restrict the study area (\SI{74.5}{\kilo\metre\squared}) to moorland peat using a mask on the \SI{10}{\metre} Sentinel-1 grid.
A pixel is kept if it lies on peat in the England Peat Map \cite{nePeatMap} and is not flagged there as improved grassland, afforested land, woodland or a linear feature, is not tree cover, open water or built-up land in ESA WorldCover 2021 \cite{zanaga2022}, and is not radar-dark water (temporal-mean VV below \SI{-18}{\deci\bel} in every orbit).
Trees and water are buffered by \SI{20}{\metre}.
The mask keeps \SI{38.8}{\kilo\metre\squared}, or 52\% of the original bounding box area (Fig.~\ref{fig:trendmap}).

From the Peatland Restoration Register we chose the area east of Crowdy Reservoir (\ang{50.62}\,N, \ang{4.62}\,W), where restoration began in March 2023.
The restoration patch is the part of the Crowdy polygon inside the mask (\SI{47.5}{\hectare}).
We also identified a contiguous control area of the same size, about \SI{2}{\kilo\metre} to the south-east, which is at least \SI{200}{\metre} from the Crowdy polygon, avoids all other restoration sites, and has similar elevation (mean \SI{291}{\metre}, against \SI{282}{\metre}, from the Copernicus GLO-30 DEM \cite{copdem}) and peat depth (100\% deep peat, against 95\%).
The control was selected using only these properties.

\subsection{Trend estimation}
\label{sec:trend}
Sentinel-1 backscatter over blanket bog has been observed to vary by about 3--6\,dB within a year \cite{toca2023potential}, which is much larger than the trends we expect to find, so we fit the trend and an annual cycle together, modelling each series $\sigma^0(t_i)$, in dB, as
\begin{equation}
  \sigma^0(t_i) = \underbrace{\beta_0 + \beta_1 (t_i - t_0)}_{\text{level + trend}}
    + \underbrace{\gamma_1 \sin 2\pi t_i + \gamma_2 \cos 2\pi t_i}_{\text{annual cycle}}
    + \varepsilon_i,
  \label{eq:model}
\end{equation}
where $t_i$ is in decimal years and $t_0 = 2024.0$.
The coefficients are estimated by ordinary least squares, and $\beta_1$ gives the trend in \si{\deci\bel\per\year}.
Consecutive acquisitions are not independent, since a wet or dry spell may last over several acquisitions, so we widen the standard error of $\beta_1$ using the effective sample size $n_\text{eff} = n(1-\rho)/(1+\rho)$, where $\rho$ is the lag-1 autocorrelation of the residuals.

We call the median VV over all masked pixels on each date the moor-wide signal, $R(t_i)$.
In order to see where on the moor the backscatter changes, we average the masked pixels into \SI{100}{\metre} cells to reduce speckle \cite{nasaSAR}, subtract $R(t_i)$ from every cell on each date, and fit Eq.~\ref{eq:model} to the result.
We call the slope from this fit the \emph{local trend}.
We take a local trend to be significant if $|\beta_1|$ is greater than two standard errors, and call a cell \emph{robust} if its local trend is significant, with the same sign, in all three orbits.
We compare the number of robust cells with the number expected if the orbits were independent (the number of cells multiplied by the product of the three orbits' fractions of significant positive trends, plus the same for negative trends); since the orbits share the same weather, we use this comparison as a guide only.

For Fig.~\ref{fig:timeseries} we average the pixels of each patch on each date, and fit Eq.~\ref{eq:model} to each patch and to the difference between them (restoration $-$ control), which removes the effect of weather common to both.

% ---------------------------------------------------------------------------
\section{Results}
% ---------------------------------------------------------------------------

\begin{figure}[t]
  \centering
  \figorbox{timeseries.pdf}{5cm}
  \caption{VV backscatter $\sigma^0$ (orbit 103) for the Crowdy restoration
  patch and its matched control. Top: patch mean on each date (points) and the
  fitted trend and annual cycle (Eq.~\ref{eq:model}; lines), with trends
  $\pm$ one standard error. Bottom: restoration $-$ control. Dashed line: start
  of restoration works; shading: June 2023, the UK's warmest June on record,
  with 68\% of average rainfall \cite{metoffice}.}
  \label{fig:timeseries}
\end{figure}

\begin{figure}[t]
  \centering
  \figorbox{trend_map.pdf}{6cm}
  \caption{Local VV trend (orbit 103) in each \SI{100}{\metre} cell, relative
  to the moor-wide signal. Black outlines: robust cells. Magenta: Crowdy
  restoration patch; gold: control (Fig.~\ref{fig:timeseries}); dashed: other
  restoration polygons. Grey: outside the peat mask.}
  \label{fig:trendmap}
\end{figure}

\subsection{Time series}
Both patches follow the annual cycle of the moor as a whole, with the highest backscatter in mid-February and the lowest in summer (Fig.~\ref{fig:timeseries}).
The largest departure from this cycle occurs during the drought of June 2023, when the mean backscatter of the control patch over June and July fell \SI{3.4}{\deci\bel} below its mean from mid-March to April, compared with \SI{2.1}{\deci\bel} for the restoration patch.
The trend in the difference between the patches is \SI{-0.34 \pm 0.23}{\deci\bel\per\year}.
However, the difference does not change steadily; it rises to \SI{+1.7}{\deci\bel} during the drought and falls to about \SI{-1}{\deci\bel} from February to April 2024, during a wet late winter and spring \cite{metoffice}.
The restoration patch also has a smaller annual cycle (amplitude $\sim$\SI{0.8}{\deci\bel}, against $\sim$\SI{1.2}{\deci\bel}).
Orbits 1 and 52 give similar results, with difference trends of \SI{-0.3}{} and \SI{-0.5}{\deci\bel\per\year}.

\subsection{Spatial pattern of trends}
Over the moor as a whole, VV increased by \SI{+0.62 \pm 0.33}{\deci\bel\per\year}.
However, the orbits do not agree on this value ($\sim$\SI{0.5}{\deci\bel\per\year} for orbit 52 and $\sim 0$ for orbit 1), and when the annual cycle is removed the increase appears to be mainly the recovery after the drop in June 2023.
Local trends are weak (Fig.~\ref{fig:trendmap}); 90\% of the 3801 cells lie within $\pm$\SI{0.5}{\deci\bel\per\year}, compared with a median standard error of $\sim$\SI{0.2}{\deci\bel\per\year}.
In total, 129 cells (3.4\%) are robust, compared with about 7 expected if the orbits were independent, and these form a few clusters, one of which lies on the eastern edge of the Crowdy restoration patch.

% ---------------------------------------------------------------------------
\section{Discussion}
% ---------------------------------------------------------------------------

\subsection{Main pattern}
We do not find a trend in backscatter over 2023--2024 which can be attributed to a change in the peat.
Local trends are weak, and the clusters of robust cells away from Crowdy should be checked against aerial imagery (for example, for felling of plantations) before they are interpreted.
At Crowdy, the evidence comes mainly from the comparison of the two patches (Fig.~\ref{fig:timeseries}), since only 2 of the 48 cells in the restoration patch are robust.
The smaller response of the restoration patch may indicate a more stable water table following ditch blocking.
Near-natural blanket bog has been found to have a smaller annual range of backscatter than previously drained bog \cite{toca2023potential}, and in Scotland, sites with blocked drains showed greater resistance to the 2018 drought than nearby near-natural sites in InSAR measurements of surface motion \cite{marshall2021naturescot}.
The lower backscatter in the restoration patch in spring 2024 could be due to water ponding behind blocked ditches, which lowers backscatter \cite{asmuss2019potential}, although a change in the vegetation layer could have a similar effect \cite{hrysiewicz2023sar}.

\subsection{Factors affecting the apparent trend}
\label{sec:factor}
The main factor affecting the apparent trend is the timing of the June 2023 drought, which occurs early in a two-year record.
A large drop near the start of the record followed by a wet year produces a positive moor-wide trend without any lasting change in the peat, and in the same way the smaller drop at Crowdy produces a negative trend in the difference between the patches (Fig.~\ref{fig:timeseries}).
The authors of \cite{marshall2021naturescot} note that short records of peat surface motion are also ``prone to influence of extreme weather events'', and suggest a baseline of about five years.

\subsection{Field observation for interpretation}
Measurements of water-table depth in both patches, taken at around the 12-day satellite interval, would show whether the smaller annual cycle in backscatter at Crowdy corresponds to a smaller range in water-table depth.
The South West Peatland Partnership already monitors water levels on Bodmin Moor \cite{swpp}, so some of these data may already be available.

\subsection{Open questions}
With only five acquisitions before the works began, we cannot determine whether the smaller response at Crowdy is due to the restoration or to an existing property of the site, which is next to Crowdy Reservoir.
Extending the record back to 2015 would show whether the restoration patch also responded less than the control during the dry summers of 2018 and 2022 \cite{ea2022}.
We would also like to extend the analysis to the whole moor, in zones of similar elevation and peat type, and to include daily rainfall data, which may help to separate the effects of wetness and vegetation, together with a comparison of VV and VH.

% ---------------------------------------------------------------------------
\section*{Code and data availability}
% ---------------------------------------------------------------------------
Code and figures: \url{https://github.com/DrPrettyman/OSC_peatland}.
Sentinel-1 data: Copernicus Sentinel data 2023--2024 \cite{torres2012}.

% ---------------------------------------------------------------------------
\begingroup
\footnotesize
\renewcommand{\section}[2]{\subsection*{References}}
\begin{thebibliography}{9}
  \setlength{\itemsep}{0pt}

  \bibitem{torres2012}
  R.~Torres et al., ``GMES Sentinel-1 mission,''
  \emph{Remote Sensing of Environment}, vol.~120, pp.~9--24, 2012.

  \bibitem{gorelick2017}
  N.~Gorelick, M.~Hancher, M.~Dixon, S.~Ilyushchenko, D.~Thau and R.~Moore,
  ``Google Earth Engine: Planetary-scale geospatial analysis for everyone,''
  \emph{Remote Sensing of Environment}, vol.~202, pp.~18--27, 2017.

  \bibitem{toca2023potential}
  L.~Toca, R.~R.~E.~Artz, C.~Smart, T.~Quaife, K.~Morrison, A.~Gimona,
  R.~Hughes, M.~H.~Hancock and D.~Klein, ``Potential for peatland water table
  depth monitoring using Sentinel-1 SAR backscatter: case study of Forsinard
  Flows, Scotland, UK,'' \emph{Remote Sensing}, vol.~15, no.~7, 1900, 2023.

  \bibitem{asmuss2019potential}
  T.~Asmu{\ss}, M.~Bechtold and B.~Tiemeyer, ``On the potential of Sentinel-1
  for high resolution monitoring of water table dynamics in grasslands on
  organic soils,'' \emph{Remote Sensing}, vol.~11, no.~14, 1659, 2019.

  \bibitem{hrysiewicz2023sar}
  A.~Hrysiewicz, E.~P.~Holohan, S.~Donohue and H.~Cushnan, ``SAR and InSAR data
  linked to soil moisture changes on a temperate raised peatland subjected to a
  wildfire,'' \emph{Remote Sensing of Environment}, vol.~291, 113516, 2023.


  \bibitem{nePeatMap}
  Natural England, ``England Peat Map'': \emph{England Peat Status, GHG and C
  Storage} (\copyright{} BGS \& NSRI) and \emph{Peatland Restoration Register},
  Natural England open data (ArcGIS feature services),
  \url{https://services.arcgis.com/JJzESW51TqeY9uat/arcgis/rest/services},
  accessed 25 Nov 2026.

  \bibitem{zanaga2022}
  D.~Zanaga et al., ``ESA WorldCover 10\,m 2021 v200,'' Zenodo, 2022,
  doi:\href{https://doi.org/10.5281/zenodo.7254221}{10.5281/zenodo.7254221}.

  \bibitem{copdem}
  European Space Agency, ``Copernicus Global Digital Elevation Model
  (GLO-30),'' doi:\href{https://doi.org/10.5270/ESA-c5d3d65}{10.5270/ESA-c5d3d65};
  accessed via Earth Engine (\texttt{COPERNICUS/DEM/GLO30\_2024\_1}).

  \bibitem{metoffice}
  Met Office, ``UK climate summaries,''
  \url{https://www.metoffice.gov.uk/research/climate/maps-and-data/summaries/index},
  accessed 25 Nov 2026.

  \bibitem{marshall2021naturescot}
  C.~Marshall, A.~V.~Bradley, R.~Andersen and D.~J.~Large, ``Using peatland
  surface motion (bog breathing) to monitor Peatland ACTION sites,''
  NatureScot Research Report 1269, 2021.

  \bibitem{bradley2022identification}
  A.~V.~Bradley, R.~Andersen, C.~Marshall, A.~Sowter and D.~J.~Large,
  ``Identification of typical ecohydrological behaviours using InSAR allows
  landscape-scale mapping of peatland condition,'' \emph{Earth Surface
  Dynamics}, vol.~10, no.~2, pp.~261--277, 2022.

  \bibitem{swpp}
  South West Peatland Partnership, ``Where we work,''
  \url{https://southwestpeatlandpartnership.co.uk/where-we-work}, accessed
  25 Nov 2026.

  \bibitem{ea2022}
  Environment Agency, ``Environment Agency chairs National Drought Group as
  parts of country move into drought,'' press release, 12 Aug 2022,
  \url{https://www.gov.uk/government/news/environment-agency-chairs-national-drought-group-as-parts-of-country-move-into-drought}.

  \bibitem{nasaSAR}
  NASA Earthdata, ``What is Synthetic Aperture Radar?,''
  \url{https://www.earthdata.nasa.gov/learn/earth-observation-data-basics/sar},
  accessed 25 Nov 2026.

  \bibitem{eeS1grd}
  Google Earth Engine Data Catalog, ``Sentinel-1 SAR GRD: C-band Synthetic
  Aperture Radar Ground Range Detected, log scaling'' (\texttt{COPERNICUS/S1\_GRD}),
  \url{https://developers.google.com/earth-engine/datasets/catalog/COPERNICUS_S1_GRD},
  accessed 25 Nov 2026.

\end{thebibliography}
\endgroup

\end{document}
